% year: תש"פ
% type: מבחן
% semester: א
% moed: ב
% part: א
% number: 1
% points: 20
% categories: func-analysis, derivatives, func-limits
% source: exams/מבחנים/תשפ/מועד ב תשפ.pdf

\begin{question}
חקרו את הפונקציה $y=\frac{x^3-1}{4x^2}$ לפי הסעיפים הבאים:
\begin{enumerate}
  \item תחום הגדרה, נקודות חיתוך עם הצירים, זוגיות, אי-זוגיות
  \item רציפות, נקודות אי-רציפות
  \item אסימפטוטות
  \item תחומי עליה וירידה, נקודות קיצון
  \item תחומי קמירות וקעירות, נקודות פיתול
  \item תיאור גרפי.
\end{enumerate}
\end{question}

\begin{hint}
כדאי לכתוב את הפונקציה בצורה $y=\frac{x}{4}-\frac{1}{4x^2}$: ממנה רואים מיד את האסימפטוטה המשופעת, והגזירה פשוטה.
\end{hint}

\begin{solution}
נסמן $f(x)=\frac{x^3-1}{4x^2}=\frac{x}{4}-\frac{1}{4x^2}$.
\begin{enumerate}
  \item \textbf{תחום הגדרה:} המכנה מתאפס רק ב-$x=0$, ולכן $D=\{x\in\R : x\neq 0\}$.

  \textbf{חיתוך עם הצירים:} $x=0$ אינו בתחום, ולכן אין חיתוך עם ציר $y$. חיתוך עם ציר $x$: $x^3-1=0 \iff x=1$, כלומר הנקודה $(1,0)$.

  \textbf{זוגיות:} $f(-x)=\frac{-x^3-1}{4x^2}$. למשל $f(1)=0$ ואילו $f(-1)=-\frac12$, כך ש-$f(-1)\neq f(1)$ וגם $f(-1)\neq -f(1)$. לכן הפונקציה אינה זוגית ואינה אי-זוגית.

  \item $f$ היא מנה של פולינומים, ולכן רציפה בכל תחום הגדרתה. הנקודה היחידה שבה $f$ אינה מוגדרת היא $x=0$. שם המונה שואף ל-$-1$ והמכנה $4x^2\to 0^+$, ולכן
  \[ \lim_{x\to 0^-} f(x)=\lim_{x\to 0^+} f(x)=-\infty , \]
  כלומר ב-$x=0$ יש נקודת אי-רציפות מסוג שני (אינסופית).

  \item \textbf{אסימפטוטה אנכית:} מהסעיף הקודם, $x=0$ היא אסימפטוטה אנכית (מ-2 הצדדים הפונקציה שואפת ל-$-\infty$). אין אסימפטוטות אנכיות נוספות כי $f$ רציפה בכל נקודה אחרת.

  \textbf{אסימפטוטה משופעת:}
  \[ m=\lim_{x\to\pm\infty}\frac{f(x)}{x}=\lim_{x\to\pm\infty}\frac{x^3-1}{4x^3}=\frac14,\qquad
     n=\lim_{x\to\pm\infty}\Big(f(x)-\frac{x}{4}\Big)=\lim_{x\to\pm\infty}\Big(-\frac{1}{4x^2}\Big)=0 . \]
  לכן $y=\frac{x}{4}$ היא אסימפטוטה משופעת גם ב-$+\infty$ וגם ב-$-\infty$. מאחר ש-$f(x)-\frac{x}{4}=-\frac{1}{4x^2}<0$, הגרף נמצא מתחת לאסימפטוטה. אין אסימפטוטה אופקית (כי $f(x)\to\pm\infty$ כאשר $x\to\pm\infty$).

  \item גוזרים:
  \[ f'(x)=\frac14+\frac{1}{2x^3}=\frac{x^3+2}{4x^3}. \]
  $f'(x)=0 \iff x^3=-2 \iff x=-\sqrt[3]{2}$. טבלת סימנים (המונה $x^3+2$ חיובי עבור $x>-\sqrt[3]{2}$, המכנה $4x^3$ חיובי עבור $x>0$):
  \begin{itemize}
    \item $x<-\sqrt[3]{2}$: מונה שלילי, מכנה שלילי, $f'>0$ --- $f$ עולה.
    \item $-\sqrt[3]{2}<x<0$: מונה חיובי, מכנה שלילי, $f'<0$ --- $f$ יורדת.
    \item $x>0$: מונה חיובי, מכנה חיובי, $f'>0$ --- $f$ עולה.
  \end{itemize}
  לכן $f$ עולה ב-$(-\infty,-\sqrt[3]{2})$ וב-$(0,\infty)$, ויורדת ב-$(-\sqrt[3]{2},0)$.
  ב-$x=-\sqrt[3]{2}$ הנגזרת עוברת מחיובית לשלילית, ולכן זו נקודת מקסימום מקומי:
  \[ f(-\sqrt[3]{2})=\frac{-2-1}{4\sqrt[3]{4}}=-\frac{3}{4\sqrt[3]{4}}\approx -0.47 . \]
  אין נקודות קיצון נוספות (ב-$x=0$ הפונקציה אינה מוגדרת).

  \item
  \[ f''(x)=\Big(\frac14+\frac12x^{-3}\Big)'=-\frac{3}{2x^4}. \]
  לכל $x\neq 0$ מתקיים $f''(x)<0$. לכן $f$ קעורה ($\cap$) בכל אחד מהקטעים $(-\infty,0)$ ו-$(0,\infty)$, ואינה קמורה ($\cup$) באף קטע. הסימן של $f''$ לא מתחלף בתוך תחום ההגדרה, ולכן אין נקודות פיתול.

  \item \textbf{תיאור גרפי:} בצד שמאל הגרף נמצא מתחת לאסימפטוטה $y=\frac{x}{4}$ וצמוד אליה כש-$x\to-\infty$. הוא עולה עד למקסימום $\big(-\sqrt[3]{2},-\frac{3}{4\sqrt[3]{4}}\big)\approx(-1.26,-0.47)$, ואז יורד ל-$-\infty$ כש-$x\to0^-$. בצד ימין הוא עולה מ-$-\infty$ (כש-$x\to0^+$), חותך את ציר $x$ ב-$(1,0)$, ומתקרב מלמטה לאסימפטוטה $y=\frac{x}{4}$ כש-$x\to\infty$. הגרף קעור בכל מקום.
\end{enumerate}
\end{solution}
